% Generated draft: requires mathematical and source review.
\documentclass[11pt]{article}
\usepackage[margin=25mm]{geometry}
\usepackage{fontspec}
\setmainfont{Latin Modern Roman}
\usepackage{amsmath,amssymb,mathtools,unicode-math}
\setmathfont{Latin Modern Math}
\usepackage{enumitem,needspace,xurl,hyperref}
\setlength{\parindent}{0pt}
\setlength{\parskip}{5pt}
\setlength{\emergencystretch}{3em}
\newcommand{\N}{\mathbb{N}}
\newcommand{\Z}{\mathbb{Z}}
\newcommand{\Q}{\mathbb{Q}}
\newcommand{\R}{\mathbb{R}}
\newcommand{\C}{\mathbb{C}}
\newcommand{\F}{\mathbb{F}}
\newcommand{\A}{\mathbb{A}}
\newcommand{\PP}{\mathbb{P}}
\newcommand{\E}{\mathbb{E}}
\newcommand{\eps}{\varepsilon}
\newcommand{\dd}{\,\mathrm{d}}
\begin{document}
\section*{UnsolvedMath ID: 20000735}
\textbf{Source problem code:} catalogue entry
\subsection{Definitions and mathematical statement}
The source asks for a dictionary between Bailey-pair methods, Bloch groups and orthogonal polynomials, possibly useful for torsion computations. In a standard convention a Bailey pair relative to \(a\) satisfies
\[
 \beta_n=\sum_{j=0}^n\frac{\alpha_j}{(q;q)_{n-j}(aq;q)_{n+j}},
 \qquad (b;q)_n=\prod_{j=0}^{n-1}(1-bq^j).
\]
The source does not choose a Bailey family, field or Bloch-group convention. For a field \(K\), take the ordinary Bloch group \(\mathcal B(K)\) as the kernel of the map \([x]\mapsto x\wedge(1-x)\) from the pre-Bloch group (symbols modulo five-term relations) to \(K^*\wedge K^*\). Torsion means a class annihilated by a nonzero integer; this does not specify its exact order.

The supplied partial dictionary concerns the standard rank-\(r\) Andrews--Gordon Bailey-chain family, \(r\ge1\), whose Nahm quadratic matrix is
\[
 A_{ij}=2\min(i,j)\qquad(1\le i,j\le r).
\]
The corresponding sums have the form \(\sum_{n\in\Z_{\ge0}^r}q^{n^{\mathsf T}An/2+B^{\mathsf T}n+C}/\prod_i(q;q)_{n_i}\), with terminal choices affecting the lower terms. The report gives the unique positive saddle in \((0,1)^r\), satisfying \(1-x_i=\prod_jx_j^{A_{ij}}\), as
\[
 c=\cos\frac\pi{2r+3},\qquad
 x_i=\frac{U_{i-1}(c)U_{i+1}(c)}{U_i(c)^2}.
\]
Here \(U_0(t)=1\), \(U_1(t)=2t\), \(U_{j+1}(t)=2tU_j(t)-U_{j-1}(t)\) are Chebyshev polynomials of the second kind. The reported distinguished class
\[
 \Xi_r=\sum_{i=1}^r[x_i]\in\mathcal B(\Q(c))
\]
is torsion over the totally real field \(\Q(c)\). No order or explicit annihilating integer is supplied.

For fixed \(r\), the report asserts that the \(r+1\) Andrews--Gordon terminal companions have distinct linear terms and product residue classes but the same saddle and same ordinary Bloch class \(\Xi_r\). Thus a dictionary whose only output is this ordinary class necessarily loses terminal-choice information. It does not identify the companions as the same function or rule out a richer dictionary retaining linear terms or additional data. Their individual linear vectors, product formulas and a general orthogonal-polynomial/Bloch correspondence are absent. The explicit Chebyshev coordinates provide one special-family link, rather than the full dictionary sought by the source.
\subsection{Short English statement}
Develop a Bailey-pair, orthogonal-polynomial and Bloch-group dictionary. The supplied report gives Chebyshev coordinates for a torsion Bloch class in an Andrews--Gordon family and shows that this ordinary class forgets which terminal companion was chosen. It does not compute all torsion orders or provide a general dictionary.
\subsection{Sources}
\begin{itemize}
\item Local UnsolvedMath registry, permanent ID \(20000735\), source code \texttt{catalogue entry}; AIM analytic number theory corpus, \texttt{aim-analytic-number-theory-notes.json}, notes[98]. Attribution: Ulam AI / UnsolvedMath Contributors, CC BY 4.0.
\item AIM \emph{Mock modular forms in combinatorics and arithmetic geometry}, Problem 2.4: \url{https://aimath.org/WWN/mockmodular/mockmodular.pdf}. The literal text repeats ``might'', without changing the mathematical request. Original remarks and literature fields are empty.
\item The summary names the Andrews--Gordon Bailey-chain family and Chebyshev/Bloch structures but supplies no complete primary references, terminal formulas or torsion orders.
\item Registry dated August 14, 2026: statement exact, status \texttt{partially\_solved}, classification \texttt{PARTIAL-PROGRESS}; attempt 1 labelled partial result, rigor proved, novelty confidence low. The explicit saddle and terminal-choice non-injectivity are machine-generated research aids, not peer reviewed, requiring independent verification before citation.
\end{itemize}
\end{document}
